\documentclass[a4,11pt]{aleph-notas}

% -- Paquetes adicionales
\usepackage{enumitem}
\usepackage{aleph-comandos}


% -- Datos  
\institucion{Facultad de Ciencias Exactas, Naturales y Ambientales}
\carrera{Catálogo STEM}
\asignatura{Álgebra lineal}
\tema{Python no. 1: Sistemas de ecuaciones}
\autor{Andrés Merino}
\fecha{Periodo 2026-1}

\logouno[0.16\textwidth]{Logos/logoPUCE_04_ac}
\definecolor{colortext}{HTML}{1957d1}
\definecolor{colordef}{HTML}{1957d1}
\fuente{montserrat}


% -- Comandos adicionales
\setlist[enumerate]{label=\roman*.}
\input{formato_ipynb}

\begin{document}

\encabezado


%%%%%%%%%%%%%%%%%%%%%%%%%%%%%%%%%%%%%%%%
\section{Sistemas de ecuaciones}
%%%%%%%%%%%%%%%%%%%%%%%%%%%%%%%%%%%%%%%%

 Podemos generar una matriz ampliada de la siguiente manera:
\begin{pycodigo}
    \begin{ipynbcodigo}\begin{lstlisting}[language=Python]
# Defino las matrices
A = Matrix([[1, 2, 3], [3, 2, 1], [1, -1, 0]])
print("A:")
display(A)
B = Matrix([[1, 0, 2], [2, 1, 0], [0, 1, 0]])
print("B:")
display(B)
# Genero la matriz ampliada
print("(A|B):")
Matrix(BlockMatrix([A, B]))
    \end{lstlisting}\end{ipynbcodigo}
    %
    \begin{ipynbsalida}
    \begin{Verbatim}
A:
    \end{Verbatim}
$\displaystyle \left[\begin{matrix}1 & 2 & 3\\3 & 2 & 1\\1 & -1 & 0\end{matrix}\right]$
    \begin{Verbatim}
B:
    \end{Verbatim}
$\displaystyle \left[\begin{matrix}1 & 0 & 2\\2 & 1 & 0\\0 & 1 & 0\end{matrix}\right]$

    \begin{Verbatim}
(A|B):
    \end{Verbatim}
$\displaystyle \left[\begin{matrix}1 & 2 & 3 & 1 & 0 & 2\\3 & 2 & 1 & 2 & 1 & 0\\1 & -1 & 0 & 0 & 1 & 0\end{matrix}\right]$
    \end{ipynbsalida}
\end{pycodigo}

%%%%%%%%%%%%%%%%%%%%%%%%%%%%%%%%%%%%%%%%
\subsection{Operaciones por filas}

\begin{itemize}
\item
    Intercambio de filas:
\begin{pycodigo}
    \begin{ipynbcodigo}\begin{lstlisting}[language=Python]
# Imprimo la matriz original
print("A:")
display(A)
# Operación de filas
print("Intercambio la fila 0 con la 2:")
A.elementary_row_op('n<->m', row1=0, row2=2)
    \end{lstlisting}\end{ipynbcodigo}
    %
    \begin{ipynbsalida}
    \begin{Verbatim}
A:
    \end{Verbatim}
    $\displaystyle \left[\begin{matrix}1 & 2 & 3\\3 & 2 & 1\\1 & -1 & 0\end{matrix}\right]$
    \begin{Verbatim}
Intercambio la fila 0 con la 2:
    \end{Verbatim}
    $\displaystyle \left[\begin{matrix}1 & -1 & 0\\3 & 2 & 1\\1 & 2 & 3\end{matrix}\right]$
    \end{ipynbsalida}
\end{pycodigo}

\item
    Multiplicar una fila por un escalar:
\begin{pycodigo}
    \begin{ipynbcodigo}\begin{lstlisting}[language=Python]
# Imprimo la matriz original
print("A:")
display(A)
# Operación de filas
print("Multiplico por 3 la fila 2:")
A.elementary_row_op('n->kn', k=3, row=2)
    \end{lstlisting}\end{ipynbcodigo}
    %
    \begin{ipynbsalida}
    \begin{Verbatim}
A:
    \end{Verbatim}
    $\displaystyle \left[\begin{matrix}1 & 2 & 3\\3 & 2 & 1\\1 & -1 & 0\end{matrix}\right]$
    \begin{Verbatim}
Multiplico por 3 la fila 2:
    \end{Verbatim}
    $\displaystyle \left[\begin{matrix}1 & 2 & 3\\3 & 2 & 1\\3 & -3 & 0\end{matrix}\right]$
    \end{ipynbsalida}
\end{pycodigo}

\item
    Sumar un múltiplo de una fila con otra:
\begin{pycodigo}
    \begin{ipynbcodigo}\begin{lstlisting}[language=Python]
# Imprimo la matriz original
print("A:")
display(A)
# Operación de filas
print("Multiplico por 3 la fila 0 y la sumo a la fila 2:")
A.elementary_row_op('n->n+km', k=3, row1=2, row2=0)
    \end{lstlisting}\end{ipynbcodigo}
    %
    \begin{ipynbsalida}
    \begin{Verbatim}
A:
    \end{Verbatim}
    $\displaystyle \left[\begin{matrix}1 & -2 & 1\\2 & 0 & -4\\3 & 1 & -2\end{matrix}\right]$
    \begin{Verbatim}
Multiplico por 3 la fila 0 y la sumo a la fila 2:
    \end{Verbatim}
    $\displaystyle \left[\begin{matrix}1 & -2 & 1\\2 & 0 & -4\\6 & -5 & 1\end{matrix}\right]$
    \end{ipynbsalida}
\end{pycodigo}

\end{itemize}

%%%%%%%%%%%%%%%%%%%%%%%%%%%%%%%%%%%%%%%%
\subsubsection{Ejemplo de Gauss-Jordan}

En caso de cometer algún error en uno de los pasos, deberá compilarse todas las celdas desde este punto.

\begin{pycodigo}
    \begin{ipynbcodigo}\begin{lstlisting}[language=Python]
# Defino la matriz
A = Matrix([[1, -2, 1], [2, 0, -4], [3, 1, -2]])
print("A:")
display(A)
# Guardo la matriz en una temporal
temp = A
    \end{lstlisting}\end{ipynbcodigo}
    %
    \begin{ipynbsalida}
    \begin{Verbatim}
A:
    \end{Verbatim}
    $\displaystyle \left[\begin{matrix}1 & -2 & 1\\2 & 0 & -4\\3 & 1 & -2\end{matrix}\right]$
    \end{ipynbsalida}
\end{pycodigo}


\begin{itemize}
\item
    Primer paso: Multiplico por \(-2\) la primera fila y lo sumo a la segunda.
\begin{pycodigo}
    \begin{ipynbcodigo}\begin{lstlisting}[language=Python]
temp = temp.elementary_row_op('n->n+km', k=-2, row1=1, row2=0)
display(temp)
    \end{lstlisting}\end{ipynbcodigo}
    %
    \begin{ipynbsalida}[2mm]
$\displaystyle \left[\begin{matrix}1 & -2 & 1\\0 & 4 & -6\\3 & 1 & -2\end{matrix}\right]$
    \end{ipynbsalida}
\end{pycodigo}

\item
    Segundo paso: Multiplico por \(-3\) la primera fila y lo sumo a la tercera.
\begin{pycodigo}
    \begin{ipynbcodigo}\begin{lstlisting}[language=Python]
temp = temp.elementary_row_op('n->n+km', k=-3, row1=2, row2=0)
display(temp)
    \end{lstlisting}\end{ipynbcodigo}
    %
    \begin{ipynbsalida}[2mm]
$\displaystyle \left[\begin{matrix}1 & -2 & 1\\0 & 4 & -6\\0 & 7 & -5\end{matrix}\right]$
    \end{ipynbsalida}
\end{pycodigo}

\item
    Tercer paso: multiplico por \(\frac{1}{4}\) la segunda fila.
\begin{pycodigo}
    \begin{ipynbcodigo}\begin{lstlisting}[language=Python]
temp = temp.elementary_row_op('n->kn', k=1/4, row=1)
display(temp)
    \end{lstlisting}\end{ipynbcodigo}
    %
    \begin{ipynbsalida}[2mm]
$\displaystyle \left[\begin{matrix}1 & -2 & 1\\0 & 1.0 & -1.5\\0 & 7 & -5\end{matrix}\right]$
    \end{ipynbsalida}
\end{pycodigo}

\item
    Cuarto paso: Multiplico por \(-7\) la segunda fila y lo sumo a la tercera.
\begin{pycodigo}
    \begin{ipynbcodigo}\begin{lstlisting}[language=Python]
temp = temp.elementary_row_op('n->n+km', k=-7, row1=2, row2=1)
display(temp)
    \end{lstlisting}\end{ipynbcodigo}
    %
    \begin{ipynbsalida}[2mm]
$\displaystyle \left[\begin{matrix}1 & -2 & 1\\0 & 1.0 & -1.5\\0 & 0 & 5.5\end{matrix}\right]$
    \end{ipynbsalida}
\end{pycodigo}

\item
    Quinto paso: Multiplico por \(\frac{1}{5.5}\) la tercera fila.
\begin{pycodigo}
    \begin{ipynbcodigo}\begin{lstlisting}[language=Python]
temp = temp.elementary_row_op('n->kn', k=1/5.5, row=2)
display(temp)
    \end{lstlisting}\end{ipynbcodigo}
    %
    \begin{ipynbsalida}[2mm]
$\displaystyle \left[\begin{matrix}1 & -2 & 1\\0 & 1.0 & -1.5\\0 & 0 & 1.0\end{matrix}\right]$
    \end{ipynbsalida}
\end{pycodigo}
\end{itemize}

%%%%%%%%%%%%%%%%%%%%%%%%%%%%%%%%%%%%%%%%
\subsubsection{Funciones para reducir una matriz}

\begin{pycodigo}
    \begin{ipynbcodigo}\begin{lstlisting}[language=Python]
# Defino la matriz
A = Matrix([[1, -2, 1], [2, 0, -4], [3, 1, -2]])
print("A:")
display(A)
    \end{lstlisting}\end{ipynbcodigo}
    %
    \begin{ipynbsalida}
    \begin{Verbatim}
A:
    \end{Verbatim}
    $\displaystyle \left[\begin{matrix}1 & -2 & 1\\2 & 0 & -4\\3 & 1 & -2\end{matrix}\right]$
    \end{ipynbsalida}
\end{pycodigo}

    
\begin{itemize}
\item
    Forma escalonada sin unos principales.
\begin{pycodigo}
    \begin{ipynbcodigo}\begin{lstlisting}[language=Python]
A.echelon_form()
    \end{lstlisting}\end{ipynbcodigo}
    %
    \begin{ipynbsalida}[2mm]
$\displaystyle \left[\begin{matrix}1 & -2 & 1\\0 & 4 & -6\\0 & 0 & 22\end{matrix}\right]$
    \end{ipynbsalida}
\end{pycodigo}

\item
    Forma escalonada reducida por fila.
\begin{pycodigo}
    \begin{ipynbcodigo}\begin{lstlisting}[language=Python]
A.rref(pivots=False)
    \end{lstlisting}\end{ipynbcodigo}
    %
    \begin{ipynbsalida}[2mm]
$\displaystyle \left[\begin{matrix}1 & 0 & 0\\0 & 1 & 0\\0 & 0 & 1\end{matrix}\right]$
    \end{ipynbsalida}
\end{pycodigo}

\item
    Rango de una matriz.
\begin{pycodigo}
    \begin{ipynbcodigo}\begin{lstlisting}[language=Python]
A.rank()
    \end{lstlisting}\end{ipynbcodigo}
    %
    \begin{ipynbsalida}
\begin{Verbatim}
3
\end{Verbatim}
    \end{ipynbsalida}
\end{pycodigo}
\end{itemize}

    

%%%%%%%%%%%%%%%%%%%%%%%%%%%%%%%%%%%%%%%%
\section{Sistemas de ecuaciones}
%%%%%%%%%%%%%%%%%%%%%%%%%%%%%%%%%%%%%%%%

 Podemos generar una matriz ampliada de la siguiente manera:
\begin{pycodigo}
    \begin{ipynbcodigo}\begin{lstlisting}[language=Python]
# Defino las matrices
A = Matrix([[1, 2, 3], [3, 2, 1], [1, -1, 0]])
print("A:")
display(A)
B = Matrix([[1, 0, 2], [2, 1, 0], [0, 1, 0]])
print("B:")
display(B)
# Genero la matriz ampliada
print("(A|B):")
Matrix(BlockMatrix([A, B]))
    \end{lstlisting}\end{ipynbcodigo}
    %
    \begin{ipynbsalida}
    \begin{Verbatim}
A:
    \end{Verbatim}
$\displaystyle \left[\begin{matrix}1 & 2 & 3\\3 & 2 & 1\\1 & -1 & 0\end{matrix}\right]$
    \begin{Verbatim}
B:
    \end{Verbatim}
$\displaystyle \left[\begin{matrix}1 & 0 & 2\\2 & 1 & 0\\0 & 1 & 0\end{matrix}\right]$

    \begin{Verbatim}
(A|B):
    \end{Verbatim}
$\displaystyle \left[\begin{matrix}1 & 2 & 3 & 1 & 0 & 2\\3 & 2 & 1 & 2 & 1 & 0\\1 & -1 & 0 & 0 & 1 & 0\end{matrix}\right]$
    \end{ipynbsalida}
\end{pycodigo}

%%%%%%%%%%%%%%%%%%%%%%%%%%%%%%%%%%%%%%%%
\subsection{Operaciones por filas}

\begin{itemize}
\item
    Intercambio de filas:
\begin{pycodigo}
    \begin{ipynbcodigo}\begin{lstlisting}[language=Python]
# Imprimo la matriz original
print("A:")
display(A)
# Operación de filas
print("Intercambio la fila 0 con la 2:")
A.elementary_row_op('n<->m', row1=0, row2=2)
    \end{lstlisting}\end{ipynbcodigo}
    %
    \begin{ipynbsalida}
    \begin{Verbatim}
A:
    \end{Verbatim}
    $\displaystyle \left[\begin{matrix}1 & 2 & 3\\3 & 2 & 1\\1 & -1 & 0\end{matrix}\right]$
    \begin{Verbatim}
Intercambio la fila 0 con la 2:
    \end{Verbatim}
    $\displaystyle \left[\begin{matrix}1 & -1 & 0\\3 & 2 & 1\\1 & 2 & 3\end{matrix}\right]$
    \end{ipynbsalida}
\end{pycodigo}

\item
    Multiplicar una fila por un escalar:
\begin{pycodigo}
    \begin{ipynbcodigo}\begin{lstlisting}[language=Python]
# Imprimo la matriz original
print("A:")
display(A)
# Operación de filas
print("Multiplico por 3 la fila 2:")
A.elementary_row_op('n->kn', k=3, row=2)
    \end{lstlisting}\end{ipynbcodigo}
    %
    \begin{ipynbsalida}
    \begin{Verbatim}
A:
    \end{Verbatim}
    $\displaystyle \left[\begin{matrix}1 & 2 & 3\\3 & 2 & 1\\1 & -1 & 0\end{matrix}\right]$
    \begin{Verbatim}
Multiplico por 3 la fila 2:
    \end{Verbatim}
    $\displaystyle \left[\begin{matrix}1 & 2 & 3\\3 & 2 & 1\\3 & -3 & 0\end{matrix}\right]$
    \end{ipynbsalida}
\end{pycodigo}

\item
    Sumar un múltiplo de una fila con otra:
\begin{pycodigo}
    \begin{ipynbcodigo}\begin{lstlisting}[language=Python]
# Imprimo la matriz original
print("A:")
display(A)
# Operación de filas
print("Multiplico por 3 la fila 0 y la sumo a la fila 2:")
A.elementary_row_op('n->n+km', k=3, row1=2, row2=0)
    \end{lstlisting}\end{ipynbcodigo}
    %
    \begin{ipynbsalida}
    \begin{Verbatim}
A:
    \end{Verbatim}
    $\displaystyle \left[\begin{matrix}1 & -2 & 1\\2 & 0 & -4\\3 & 1 & -2\end{matrix}\right]$
    \begin{Verbatim}
Multiplico por 3 la fila 0 y la sumo a la fila 2:
    \end{Verbatim}
    $\displaystyle \left[\begin{matrix}1 & -2 & 1\\2 & 0 & -4\\6 & -5 & 1\end{matrix}\right]$
    \end{ipynbsalida}
\end{pycodigo}

\end{itemize}

%%%%%%%%%%%%%%%%%%%%%%%%%%%%%%%%%%%%%%%%
\subsubsection{Ejemplo de Gauss-Jordan}

En caso de cometer algún error en uno de los pasos, deberá compilarse todas las celdas desde este punto.

\begin{pycodigo}
    \begin{ipynbcodigo}\begin{lstlisting}[language=Python]
# Defino la matriz
A = Matrix([[1, -2, 1], [2, 0, -4], [3, 1, -2]])
print("A:")
display(A)
# Guardo la matriz en una temporal
temp = A
    \end{lstlisting}\end{ipynbcodigo}
    %
    \begin{ipynbsalida}
    \begin{Verbatim}
A:
    \end{Verbatim}
    $\displaystyle \left[\begin{matrix}1 & -2 & 1\\2 & 0 & -4\\3 & 1 & -2\end{matrix}\right]$
    \end{ipynbsalida}
\end{pycodigo}


\begin{itemize}
\item
    Primer paso: Multiplico por \(-2\) la primera fila y lo sumo a la segunda.
\begin{pycodigo}
    \begin{ipynbcodigo}\begin{lstlisting}[language=Python]
temp = temp.elementary_row_op('n->n+km', k=-2, row1=1, row2=0)
display(temp)
    \end{lstlisting}\end{ipynbcodigo}
    %
    \begin{ipynbsalida}[2mm]
$\displaystyle \left[\begin{matrix}1 & -2 & 1\\0 & 4 & -6\\3 & 1 & -2\end{matrix}\right]$
    \end{ipynbsalida}
\end{pycodigo}

\item
    Segundo paso: Multiplico por \(-3\) la primera fila y lo sumo a la tercera.
\begin{pycodigo}
    \begin{ipynbcodigo}\begin{lstlisting}[language=Python]
temp = temp.elementary_row_op('n->n+km', k=-3, row1=2, row2=0)
display(temp)
    \end{lstlisting}\end{ipynbcodigo}
    %
    \begin{ipynbsalida}[2mm]
$\displaystyle \left[\begin{matrix}1 & -2 & 1\\0 & 4 & -6\\0 & 7 & -5\end{matrix}\right]$
    \end{ipynbsalida}
\end{pycodigo}

\item
    Tercer paso: multiplico por \(\frac{1}{4}\) la segunda fila.
\begin{pycodigo}
    \begin{ipynbcodigo}\begin{lstlisting}[language=Python]
temp = temp.elementary_row_op('n->kn', k=1/4, row=1)
display(temp)
    \end{lstlisting}\end{ipynbcodigo}
    %
    \begin{ipynbsalida}[2mm]
$\displaystyle \left[\begin{matrix}1 & -2 & 1\\0 & 1.0 & -1.5\\0 & 7 & -5\end{matrix}\right]$
    \end{ipynbsalida}
\end{pycodigo}

\item
    Cuarto paso: Multiplico por \(-7\) la segunda fila y lo sumo a la tercera.
\begin{pycodigo}
    \begin{ipynbcodigo}\begin{lstlisting}[language=Python]
temp = temp.elementary_row_op('n->n+km', k=-7, row1=2, row2=1)
display(temp)
    \end{lstlisting}\end{ipynbcodigo}
    %
    \begin{ipynbsalida}[2mm]
$\displaystyle \left[\begin{matrix}1 & -2 & 1\\0 & 1.0 & -1.5\\0 & 0 & 5.5\end{matrix}\right]$
    \end{ipynbsalida}
\end{pycodigo}

\item
    Quinto paso: Multiplico por \(\frac{1}{5.5}\) la tercera fila.
\begin{pycodigo}
    \begin{ipynbcodigo}\begin{lstlisting}[language=Python]
temp = temp.elementary_row_op('n->kn', k=1/5.5, row=2)
display(temp)
    \end{lstlisting}\end{ipynbcodigo}
    %
    \begin{ipynbsalida}[2mm]
$\displaystyle \left[\begin{matrix}1 & -2 & 1\\0 & 1.0 & -1.5\\0 & 0 & 1.0\end{matrix}\right]$
    \end{ipynbsalida}
\end{pycodigo}
\end{itemize}

%%%%%%%%%%%%%%%%%%%%%%%%%%%%%%%%%%%%%%%%
\subsubsection{Funciones para reducir una matriz}

\begin{pycodigo}
    \begin{ipynbcodigo}\begin{lstlisting}[language=Python]
# Defino la matriz
A = Matrix([[1, -2, 1], [2, 0, -4], [3, 1, -2]])
print("A:")
display(A)
    \end{lstlisting}\end{ipynbcodigo}
    %
    \begin{ipynbsalida}
    \begin{Verbatim}
A:
    \end{Verbatim}
    $\displaystyle \left[\begin{matrix}1 & -2 & 1\\2 & 0 & -4\\3 & 1 & -2\end{matrix}\right]$
    \end{ipynbsalida}
\end{pycodigo}

    
\begin{itemize}
\item
    Forma escalonada sin unos principales.
\begin{pycodigo}
    \begin{ipynbcodigo}\begin{lstlisting}[language=Python]
A.echelon_form()
    \end{lstlisting}\end{ipynbcodigo}
    %
    \begin{ipynbsalida}[2mm]
$\displaystyle \left[\begin{matrix}1 & -2 & 1\\0 & 4 & -6\\0 & 0 & 22\end{matrix}\right]$
    \end{ipynbsalida}
\end{pycodigo}

\item
    Forma escalonada reducida por fila.
\begin{pycodigo}
    \begin{ipynbcodigo}\begin{lstlisting}[language=Python]
A.rref(pivots=False)
    \end{lstlisting}\end{ipynbcodigo}
    %
    \begin{ipynbsalida}[2mm]
$\displaystyle \left[\begin{matrix}1 & 0 & 0\\0 & 1 & 0\\0 & 0 & 1\end{matrix}\right]$
    \end{ipynbsalida}
\end{pycodigo}

\item
    Rango de una matriz.
\begin{pycodigo}
    \begin{ipynbcodigo}\begin{lstlisting}[language=Python]
A.rank()
    \end{lstlisting}\end{ipynbcodigo}
    %
    \begin{ipynbsalida}
\begin{Verbatim}
3
\end{Verbatim}
    \end{ipynbsalida}
\end{pycodigo}
\end{itemize}

%%%%%%%%%%%%%%%%%%%%%%%%%%%%%%%%%%%%%%%%
\subsection{Resolución de sistemas de ecuaciones}

Consideremos el sistema de ecuaciones \(Ax=b\), donde 
\[
    A = \begin{pmatrix}1 & 2 & 1\\0 & 1 & 2\\2 & 4 & 0\end{pmatrix}
    \quad\text{y}\quad
    b=\begin{pmatrix}7\\12\\4\end{pmatrix}.
\]
Podemos resolverlo con el comando \texttt{solve} de la siguiente manera:
\begin{pycodigo}
    \begin{ipynbcodigo}\begin{lstlisting}[language=Python]
# Defino las matrices
A = Matrix([[1, 2, 1], [0, 1, 2], [2, 4, 0]])
b = Matrix([7, 12, 4])
# Solución
A.solve(b)
    \end{lstlisting}\end{ipynbcodigo}
    %
    \begin{ipynbsalida}[2mm]
$\displaystyle \left[\begin{matrix}-2\\2\\5\end{matrix}\right]$
    \end{ipynbsalida}
\end{pycodigo}

Sin embargo, si el sistema no tiene solución única, nos dará un error. En estos casos, podemos utilizar \texttt{gauss\_jordan\_solve}:

\begin{pycodigo}
    \begin{ipynbcodigo}\begin{lstlisting}[language=Python]
# Solución
sol, par = A.gauss_jordan_solve(b)
sol
    \end{lstlisting}\end{ipynbcodigo}
    %
    \begin{ipynbsalida}[2mm]
$\displaystyle \left[\begin{matrix}-2\\2\\5\end{matrix}\right]$
    \end{ipynbsalida}
\end{pycodigo}

Este comando, a más de determinar la solución, presenta los parámetros que esta tendrá en el caso de no tener solución única. Para esto, consideremos el sistema de ecuaciones \(Ax=b\), con 
\[
    A=\begin{pmatrix}1 & 2 & 1 & 1\\1 & 2 & 2 & -1\\2 & 4 & 0 & 6\end{pmatrix}
    \quad\text{y}\quad
    b=\begin{pmatrix}7\\12\\4\end{pmatrix}.
\]

\begin{pycodigo}
    \begin{ipynbcodigo}\begin{lstlisting}[language=Python]
# Definimos la matrices
A = Matrix([[1, 2, 1, 1], [1, 2, 2, -1], [2, 4, 0, 6]])
b = Matrix([7, 12, 4])
# Solución
sol, par = A.gauss_jordan_solve(b)
print("Solución:")
display(sol)
print("Parámetros:")
display(par)
    \end{lstlisting}\end{ipynbcodigo}
    %
    \begin{ipynbsalida}[2mm]
    \begin{Verbatim}
Solución:
    \end{Verbatim}
    $\displaystyle \left[\begin{matrix}- 2 \tau_{0} - 3 \tau_{1} + 2\\\tau_{0}\\2 \tau_{1} + 5\\\tau_{1}\end{matrix}\right]$
    \begin{Verbatim}
Parámetros:
    \end{Verbatim}
    $\displaystyle \left[\begin{matrix}\tau_{0}\\\tau_{1}\end{matrix}\right]$
    \end{ipynbsalida}
\end{pycodigo}



    



\end{document}